\textbf{结论名称：} 球的内接正方体/长方体/正四面体的尺寸关系

\textbf{适用条件：} 正方体、长方体或正四面体的所有顶点均在球面上（即球为其外接球）。

\textbf{变量说明：} \(R\) 为球半径；正方体棱长为 \(a\)；长方体长、宽、高分别为 \(a,b,c\)；正四面体棱长为 \(a\)。

\textbf{核心公式：}
\[
\boxed{
\begin{aligned}
\text{正方体:}&\quad 2R = a\sqrt{3},\quad R = \frac{\sqrt{3}}{2}a \\[4pt]
\text{长方体:}&\quad 2R = \sqrt{a^2+b^2+c^2} \\[4pt]
\text{正四面体:}&\quad 2R = \frac{\sqrt{6}}{2}a,\quad R = \frac{\sqrt{6}}{4}a
\end{aligned}
}
\]

\textbf{使用场景：} 已知球半径求内接几何体的棱长，或已知棱长求球体积、表面积，亦可用于多面体与球的截面问题中快速转化。

\textbf{简单例子：} 球的内接正方体棱长为 \(2\)，则球半径 \(R = \frac{\sqrt{3}}{2}\times 2 = \sqrt{3}\)，球的表面积 \(S = 4\pi R^2 = 12\pi\)。